% -----------------------------------------------------------------------------
% 2. Background and Related Work — deux paragraphes : ce qui est établi (grappes
% de citations), ce qui reste ouvert (énoncé du gap). Budget : ~280 mots.
% -----------------------------------------------------------------------------
Automatic detection of potentially unfair clauses in terms of service is by now
well established: \corpus{} and its successors detect and explain unfairness in
English and across languages \cite{lippi2019claudette,lagioia2019deep,
ruggeri2022memory,drawzeski2021corpus,galassi2024multilingual,liepina2020claudettetool},
and the task is part of standard legal benchmarks \cite{chalkidis2022lexglue}.
Compact domain-specific encoders remain competitive with, and often better than,
general-purpose large models on clause classification
\cite{chalkidis2020legalbert,dominguezolmedo2025lawma,panarelli2025worth,frasheri2024llm}.
Language models are increasingly used as semantic annotators of legal text, with
documented gains in cost and documented risks of anchoring and circularity
\cite{savelka2023unreasonable,gilardi2023chatgpt,schepers2025price,
li2023coannotating,choi2024llmeffect}, and, as judges, with known divergence from
human evaluators \cite{zheng2023judging,bavaresco2025judges,thalken2023legal}.
On the measurement side, chance-corrected agreement for set-valued annotation is
a mature subject \cite{krippendorff2018content,passonneau2006masi,
artstein2008intercoder,marchal2022multilabel,gwet2008ac1}, and a meta-analysis of
annotation studies identifies the number of categories as one of the strongest
determinants of agreement \cite{bayerl2011metaanalysis}.

What is missing is a reliable, measured and published thematic layer on the
reference benchmark itself. Contract-topic resources adjacent to ours are either
single-annotated \cite{braun2022topics}, defined for a different document type
\cite{tuggener2020ledgar,hendrycks2021cuad,wilson2016opp115}, or reported without
chance-corrected agreement \cite{loffler2025chilean,braun2024agbde}; none places
human annotators and language models on the same vocabulary and the same
sentences. The consequence is methodological as much as practical: because the
multi-label format and the granularity of the taxonomy are almost never varied on
the same data, it is not known whether the agreement penalty attributed to
multi-label annotation is a property of the format or of the label set it is
applied to --- a question \cite{bayerl2011metaanalysis} raises and leaves open.
